%------------------------------------------------------------------------------%
\begin{exercises}

% 71 - Stewart pt pg 634
%------------------------------------------------------------------------------%
\begin{exercise}
  Suponha que que $(a_n)$ é uma sequência decrescente
  e que todos os termos estão entre os números 5 e 8.
  Explique por que a sequência tem um limite.
  O que você pode dizer sobre o valor do limite?
\end{exercise}

% 112 - Thomas pg 12
%------------------------------------------------------------------------------%
\begin{exercise}
Determine se a sequência
\[
  a_n = \frac{(2n+3)!}{(n+1)!}
\]
é monotônica e se é limitada.
\begin{answer}
  Para determinar que a sequência é crescente vamos verificar que $a_n<a_{n+1}$
  \begin{align*}
  	a_n                        & < a_{n+1}                             \\[2mm]
  	\frac{(2n+3)!}{(n+1)!}     & < \frac{(2(n+1)+3)!}{((n+1)+1)!}      \\[2mm]
  	\frac{(2n+3)!}{(n+1)!}     & < \frac{(2n+5)!}{(n+2)!}              \\[2mm]
  	\frac{(n+2)!}{(n+1)!}      & < \frac{(2n+5)!}{(2n+3)!}             \\[2mm]
  	\frac{(n+2)(n+1)!}{(n+1)!} & < \frac{(2n+5)(2n+4)(2n+3)!}{(2n+3)!} \\
  	(n+2)                      & < (2n+5)(2n+4)                        \\
  	n+2                        & <  4n^2 +8n + 10n +20                 \\
  	n+2                        & <  4n^2 +18n +20                      \\
  	0                          & <  4n^2 +17n +18
  \end{align*}
  Como todas esses desigualdades são equivalentes ($\Leftrightarrow$)
  concluímos que a sequência é crescente e portanto monótona.

  Para verificar se ela é limitada vamos desenvolver a expressão do termo geral
  \begin{align*}
  	a_n & = \frac{(2n+3)!}{(n+1)!}                       \\
  	    & = \frac{(2n+3)(2n+2)\cdots(n+2)(n+1)!}{(n+1)!} \\
  	    & = (2n+3)(2n+2)\cdots(n+2)
  \end{align*}
  Observamos que o termo geral não é limitado.
\end{answer}
\end{exercise}

% 111, 113, 114 - Thomas pg 12
%------------------------------------------------------------------------------%
\begin{exercise}
  Determine se cada sequência é monotônica e se é limitada
  \begin{mathlist}
    \item a_n = \frac{3n+1}{n+1}
    \item a_n = \frac{2^n\,3^n}{n!}
    \item a_n = 2 - \frac{2}{n} - \frac{1}{2^n}
  \end{mathlist}
\end{exercise}

% 72-78 Steward pt pg 634
%------------------------------------------------------------------------------%
\begin{exercise}
  Determine se a sequência dada é crescente, decrescente ou não
  monótona. Verifique se a sequência é limitada.
  \begin{mathlist}[2]
    \item a_n = (-2)^{n+1}
    \item a_n = \frac{2}{2n+5}
    \item a_n = \frac{2n-3}{3n+4}
    \item a_n = (-1)^n\,n
    \item a_n = ne^{-n}
    \item a_n = \frac{n}{n^2+1}
    \item a_n = n + \frac{1}{n}
  \end{mathlist}
\end{exercise}

% 80 Steward pt pg 634
%------------------------------------------------------------------------------%
\begin{exercise}
  Considere a sequência $(a_n)$ definida como
  \[
    a_1 = \sqrt{2}
    \qquad
    a_{n+1} = \sqrt{2+a_n}
  \]
  \vspace{-\parsep}
  \begin{alphalist}
    \item Mostre por indução que $(a_n)$ é crescente.
    \item Mostre que ela é limitada superiormente por~3.
    \item Aplique o Teorema da Sequencia Monótona para mostrar que
          $(a_n)$ é convergente.
    \item Calcule o limite de $(a_n)$.
  \end{alphalist}
\end{exercise}

% 81 Steward pt pg 634
%------------------------------------------------------------------------------%
\begin{exercise}
  Mostre que a sequência
  \[
  a_1 = 1
  \qquad
  a_{n+1} = 3 - \frac{1}{a_n}
  \]
  é crescente e $a_n<3$ para todo $n$.
  Verifique se $(a_n)$ é convergente e se for calcule seu limite.
\end{exercise}

% 82 Steward pt pg 634
%------------------------------------------------------------------------------%
\begin{exercise}
  Mostre que a sequência
  \[
    a_1 = 2
    \qquad
    a_{n+1} = \frac{1}{3-a_n}
  \]
  é decrescente e satisfaz $0<a_n\leq 3$ para todo $n$.
  Verifique se $(a_n)$ é convergente e se for calcule seu limite.
\end{exercise}

%------------------------------------------------------------------------------%
%\begin{exercise}
%\begin{answer}
%\end{answer}
%\end{exercise}

%------------------------------------------------------------------------------%
%\begin{exercise}
%\begin{mathlist}[2]
%  \item
%  \item
%  \item
%\end{mathlist}
%\begin{answer}
%  \begin{alphalist}[3]
%    \item
%    \item
%    \item
%  \end{alphalist}
%\end{answer}
%\end{exercise}

\end{exercises}
%------------------------------------------------------------------------------%
